% !TEX root = ../monografia.tex
% Usa glossaries já carregado e preserva a lista de siglas do IFSC.
% Sintaxe: \registrarsigla{chave}{SIGLA}{nome por extenso}
\newcommand{\registrarsigla}[3]{%
    \newacronym{#1}{#2}{#3}%
    \abreviatura{#2}{#3}%
}

\registrarsigla{yolo}{YOLO}{You Only Look Once}
\registrarsigla{ram}{RAM}{memória de acesso aleatório}

% Fundamentos de imagens, detecção e inferência embarcada.
\registrarsigla{rgb}{RGB}{Red, Green and Blue}
\registrarsigla{bgr}{BGR}{Blue, Green and Red}
\registrarsigla{jpeg}{JPEG}{Joint Photographic Experts Group}
\registrarsigla{cnn}{CNN}{Convolutional Neural Network}
\registrarsigla{relu}{ReLU}{Rectified Linear Unit}
\registrarsigla{iou}{IoU}{Intersection over Union}
\registrarsigla{nms}{NMS}{Non-Maximum Suppression}
\registrarsigla{onnx}{ONNX}{Open Neural Network Exchange}
\registrarsigla{cpu}{CPU}{unidade central de processamento}
\registrarsigla{simd}{SIMD}{Single Instruction, Multiple Data}
\registrarsigla{csi2}{CSI-2}{Camera Serial Interface 2}
\registrarsigla{lpddr2}{LPDDR2}{Low Power Double Data Rate 2}

% Mantidas para compatibilidade com os demais capítulos do modelo.
\registrarsigla{ifsc}{IFSC}{Instituto Federal de Santa Catarina}
\registrarsigla{daeln}{DAELN}{Departamento Acadêmico de Eletrônica}
\registrarsigla{ihm}{IHM}{Interface Homem-Máquina}
\registrarsigla{iot}{IoT}{Internet das Coisas -- \textit{Internet of Things}}
